\documentclass{article}
\usepackage{../official_template/spconf,amsmath,graphicx,booktabs}
\usepackage{url}

\title{CloneBlock: Session-Level Aggregation-Aware Defense against\\
Realistic Zero-Shot Voice Cloning}

% AUTHOR TODO: Replace every placeholder and add ORCIDs in the submission system.
\name{Anonymous Author List Pending Confirmation}
\address{Affiliations Pending Confirmation}

\begin{document}
\maketitle

\begin{abstract}
% AUTHOR TODO (100--150 words): problem, realistic protocol, method change,
% principal security result, speech-quality qualification, and XTTS finding.
% Do not present the clean-information oracle as a certified upper bound.
This submission draft intentionally leaves the abstract for author approval.
\end{abstract}

\begin{keywords}
voice cloning, adversarial defense, speaker verification, telephony, speech security
\end{keywords}

\section{Introduction}
\label{sec:introduction}
% AUTHOR TODO: Motivate multi-utterance/session-level enrollment attacks.
% Distinguish the revised contribution from the prior Interspeech submission.
% Keep the contribution list to three auditable claims:
% (1) realistic aggregation-aware protocol,
% (2) channel-aware training/evaluation and uncertainty,
% (3) cross-encoder plus small-scale end-to-end cloning evidence.

\section{Threat Model and Method}
\label{sec:method}

\subsection{Session-level aggregation-aware protocol}
% AUTHOR TODO: Define attacker knowledge, candidate budget K, aggregation rule,
% enrollment boundary, and whether the target speaker contributes any clean
% reference information. Explain why repeated utterances form one session.

Let $x$ denote clean speech, $D(x)$ the defended waveform, and
$T(\cdot;\theta)$ a sampled telephone channel. We use
$x_{\mathrm{tele}}=T(x;\theta)$ and
$x_{\mathrm{tele\_def}}=T(D(x);\theta)$. The primary protocol computes attack
selection evidence from $x_{\mathrm{tele\_def}}$; selection from $D(x)$ is an
easier pre-channel diagnostic only.

\subsection{Realistic and information-privileged objectives}
% AUTHOR TODO: Insert exact equations from src/metrics/session_attack.py after
% line-by-line author verification. Name the deployed objective
% `realistic_selected` and the diagnostic objective `reference_v0p2`.
The realistic view uses only observations available after defense and channel
processing. A separate clean-information oracle is reported as an
information-privileged reference, not as a guaranteed numerical upper bound.

\subsection{Training objective and perturbation budget}
% AUTHOR TODO: Explain the realistic-objective retraining, validation sweep,
% selected checkpoint, norm/budget convention, and causal streaming boundary.

\section{Experimental Protocol}
\label{sec:experiments}
% AUTHOR TODO: Add exact VCTK release, speaker split, utterance counts, sample
% rate, seeds, codecs, ASR model, ECAPA/x-vector model IDs, and XTTS version.

We report $K{=}16$ success and AUCK. AUCK is the normalized mean (equivalently,
the normalized discrete area) across the prespecified $K$ curve; the exact
$K$ grid must be stated in the final text. Session-bootstrap 95\% confidence
intervals use the frozen split and seed recorded in the evidence manifest.
EER denotes equal error rate, and target--impostor margin is the target-speaker
score minus the strongest impostor score.

\section{Results}
\label{sec:results}

\begin{table}[t]
  \centering
  \caption{Session-level ECAPA attack metrics with the primary
  \texttt{tele\_defended} selection source. Higher attack success and AUCK are
  worse for the defense. Confidence intervals remain in the machine-readable
  evidence table and will be added after layout selection.}
  \label{tab:main-security}
  \small
  \setlength{\tabcolsep}{4pt}
  \begin{tabular}{lcc}
    \toprule
    Condition & $K{=}16$ & AUCK \\
    \midrule
    No defense & 0.9873 & 0.8889 \\
    Clean-reference diagnostic & 0.9444 & 0.8497 \\
    Realistic selected & \textbf{0.9130} & \textbf{0.8208} \\
    \bottomrule
  \end{tabular}
\end{table}

\begin{table}[t]
  \centering
  \caption{Realistic ECAPA attack metrics under channels/codecs.}
  \label{tab:channel-robustness}
  \small
  \setlength{\tabcolsep}{5pt}
  \begin{tabular}{lcc}
    \toprule
    Channel & $K{=}16$ & AUCK \\
    \midrule
    Telephone transform only & 0.9130 & 0.8208 \\
    Opus 16 kb/s & 0.9096 & 0.8143 \\
    G.711 & 0.9111 & 0.8206 \\
    \bottomrule
  \end{tabular}
\end{table}

\begin{table*}[t]
  \centering
  \caption{Utility and end-to-end evidence. WER $\Delta$ is defended minus
  undefended under the same channel draw. STOI compares
  \texttt{tele\_clean} with \texttt{tele\_defended}. XTTS reports the paired
  x-vector target--impostor-margin change (defended minus baseline).}
  \label{tab:quality-cloning}
  \small
  \setlength{\tabcolsep}{4pt}
  \begin{tabular}{lccc}
    \toprule
    Channel / task & Estimate & 95\% CI & Auxiliary metric \\
    \midrule
    No codec WER $\Delta$ & $+0.0025$ & $[-0.0020,\ 0.0071]$ & STOI 0.9844 \\
    Opus WER $\Delta$ & $+0.0023$ & $[-0.0050,\ 0.0101]$ & STOI 0.9068 \\
    G.711 WER $\Delta$ & $+0.0047$ & $[-0.0031,\ 0.0123]$ & STOI 0.9772 \\
    XTTS margin $\Delta$, $K{=}1$ & $-0.00246$ & $[-0.00504,\ -0.00019]$ & x-vector \\
    XTTS margin $\Delta$, $K{=}16$ & $-0.00296$ & $[-0.00521,\ -0.00052]$ & x-vector \\
    \bottomrule
  \end{tabular}
\end{table*}


% AUTHOR TODO: Interpret, rather than merely repeat, the tables. Required
% qualifications: defended-source selection is diagnostic; absolute ASR WER is
% weak even on original clean speech; quality deltas cross zero; XTTS effects
% are encoder-dependent; no unaudited ClearMask/Enkidu numbers may be reused.

\section{Limitations and Conclusion}
\label{sec:conclusion}
% AUTHOR TODO: Discuss one corpus, small XTTS subset, two speaker encoders,
% missing auditable baseline reruns, and the distinction between speaker-score
% displacement and practical clone rejection. End with the narrow supported
% conclusion, not a universal prevention claim.

\clearpage
\bibliographystyle{../official_template/IEEEbib}
\bibliography{refs}

\end{document}
